\section{Proposal Context and Scope}
\label{sec:proposal-context}

The paper studies the proposal cluster that defines modern Swift concurrency behavior. \code{isolated} parameters and actor-instance control come from SE-0313~\cite{se0313}. Region-based isolation and flow-sensitive transfer come from SE-0414~\cite{se0414}. Method and key-path \code{Sendable} inference come from SE-0418~\cite{se0418}. Caller isolation inheritance and \code{\#isolation} come from SE-0420~\cite{se0420}. Explicit \code{sending} parameters and results come from SE-0430~\cite{se0430}. Dynamic \code{@isolated(any)} function values come from SE-0431~\cite{se0431}. The caller-actor execution model for nonisolated async functions is revised by SE-0461~\cite{se0461}.

The goal is not to restate every rule those proposals contain. The goal is to extract the smallest account that still explains the behaviors programmers most often find irregular:

\begin{itemize}
  \item closure conversion and the loss or preservation of inherited actor context,
  \item the difference between same-isolation binding and cross-isolation transfer for non-\code{Sendable} values,
  \item the boundary between \code{@nonisolated async} and \code{@concurrent},
  \item dynamic isolation through \code{@isolated(any)},
  \item and caller-side evidence mechanisms such as \code{isolated} parameters and \code{\#isolation}.
\end{itemize}

The resulting model is intentionally selective. It is not a full mechanization of Swift, and it does not try to formalize every parser rule, standard-library API, or SIL optimization. It instead aims for explanatory leverage: a vocabulary strong enough to make the recurring diagnostics predictable.

\section{Why the Surface Model Feels Irregular}
\label{sec:irregular}

The most common informal explanation of Swift Concurrency is that code is either isolated or nonisolated. That explanation is too coarse for at least three reasons.

First, function values and function bodies live on different semantic layers.

\codefile{snippets/handler-location.swift}

The property \code{handler} lives in the region of the enclosing declaration, while the function body that the property stores runs according to the isolation annotation on the function type. A single phrase such as ``the closure is on \code{@FooActor}'' is therefore ambiguous. One must distinguish storage location from execution capability.

Second, accessibility and capturability are not the same relation. A task-isolated value may be directly readable in some contexts and still be illegal to capture into an actor-isolated closure. This distinction is crucial for understanding nonisolated async parameters after SE-0461 and for explaining why \code{@isolated(any)} closures can capture some values that \code{@MainActor} closures cannot~\cite{se0461}.

Third, boundary crossing is not the whole story. Same-isolation ordinary calls do not consume non-\code{Sendable} values, but they may refine a disconnected value into an actor-bound region. Same-isolation \code{sending} calls preserve the original region instead of binding it. Cross-isolation calls may consume the argument, while cross-isolation results require a separate explicit \code{sending} guarantee. The model therefore needs more than a boolean boundary test; it needs a caller-side environment that changes over time.

\begin{definitionbox}{Three semantic relations}
The paper distinguishes three relations that are often conflated in surface reasoning. $\Accessible(@\kappa)$ determines which regions can be used directly in a capability~$@\kappa$. $\Capturable(@\kappa)$ determines which regions may be captured into a closure checked under~$@\kappa$. $\canSend(@\kappa, @\iota, [\sending])$ determines when a call site performs an explicit or implicit transfer.
\end{definitionbox}

These distinctions also explain why global actors and actor instances behave differently under closure inference. A global-actor closure can inherit actor context without capturing a concrete actor value. An actor-instance closure keeps instance isolation only if the isolated parameter or \code{self} is actually captured. That capture-sensitive branch is easy to miss in proposal prose, but it is essential for matching compiler behavior.

The next section turns these informal observations into a judgment form and a small set of auxiliary definitions.
